\subsection*{P2 task 5: noise-texture test (pre-specified)}
Written 2 Nov 2026, \textasciitilde{}17:15, before any NPS-matched noise was generated or scored. The SHA-256 hash of this file is recorded in \texttt{analysis/p2\_texture/PRESPEC\_SHA256.txt} before the generation step. Any later change to this plan is listed at the end, with the date and the reason.

\subsubsection*{Why}
In the draft, the claim that noise texture explains the different effects of white, correlated and projection-domain noise is argued but not tested. The three noise models differ in texture, level, spatial uniformity and insertion path all at once. For example, "80 HU" image-domain noise adds 67 HU over the lung crop, while projection noise at 10 \% dose adds a median of about 102 HU. This test holds the level and the insertion path fixed and varies only the texture.

\subsubsection*{Conditions (LUNA16 test folds 0-3, 356 scans; S3 and S4 of each fold; paired, identical inputs)}
All noise is added to the ORIGINAL full-field LUNA16 image (HU, native grid). It then goes through the unchanged preprocessing with the same lung mask, exactly the path used for the projection-domain noise (\texttt{scripts/17\_luna\_projnoise.py}). Seeds are fixed per (scan, condition). The noise is stationary, zero-mean, Gaussian and additive.

\begin{itemize}
\item \textbf{T1 \texttt{tex\_nps}}: in-plane noise power spectrum (NPS) equal to the measured NPS of the real TCIA noise (10 \%-dose minus routine image). The B50f chest pairs give the spectrum for sharp-kernel scans and the B30f abdomen pairs for soft-kernel scans, the same kernel-class split as the projection noise. Slices are independent along z, as in the projection noise. The radial NPS is measured in cycles/mm and applied at each scan's own pixel spacing. It is zero above the highest frequency measured on TCIA.
\item \textbf{T2 \texttt{tex\_gauss}}: the texture of the correlated image-domain noise of the draft: white noise low-pass filtered with a 3D Gaussian of σ = 1.41 mm in-plane and 2.5 mm along z (2 voxels on the detector grid; the "1.4 mm grain" condition), here in physical units on the native grid.
\item \textbf{T3 \texttt{tex\_white}}: white noise, independent per native voxel.
\end{itemize}

\textbf{Level:} for every scan and condition, the added-noise SD over the preprocessed lung crop (preprocessed noisy volume minus the clean archive) is matched to that scan's projection-domain noise SD at 10 \% dose (\texttt{noise\_sd\_added\_hu} in \texttt{projnoise/a0.1/.../\textless{}uid\textgreater{}.json}). The scale is set in two passes; the accepted tolerance is ±3 \% per scan, and any scan outside it is reported. T1-T3 and proj0.1 therefore have the same level, crop and insertion path. They differ in texture (T1 vs T2 vs T3) or, for T1 vs proj0.1, only in the spatial non-uniformity and streak structure of the projection noise.

\subsubsection*{Known-answer checks before scoring}
\begin{enumerate}
\item The generated T1 noise, measured with the same ROI method, reproduces the target radial NPS shape (normalised L1 distance \textless{} 0.05). The generated T2 and T3 noise have the expected flat or Gaussian shape.
\item With zero amplitude, the output is byte-identical to the clean archive (3 scans).
\item Every scan's achieved SD is within ±3 \% of its target.
\item The new inference path reproduces the existing proj0.1 candidates of one scan exactly.
\end{enumerate}

\subsubsection*{Primary family F4 (Holm within the family, α = 0.05, m = 2)}
For each model M ∈ \{S3, S4\}: contrast CPM\_M(T1) − CPM\_M(T2), on the same paired bootstrap as F1-F3 (10,000 resamples over scans, common resamples, two-sided P, floor 2/10,000). The texture hypothesis predicts \textgreater{} 0: at equal level and path, noise with the real texture harms less than Gaussian-blob noise.

\subsubsection*{Pre-specified equivalence check (secondary, not in Holm)}
For each model, CPM\_M(T1) − CPM\_M(proj0.1), 95 \% CI inside ±0.02 → "stationary noise with the real NPS reproduces the projection-domain result".

\subsubsection*{Secondary (reported, unadjusted)}
\begin{itemize}
\item Each model's drop from clean under T1, T2 and T3 (paired CIs).
\item CPM\_M(T3) − CPM\_M(T1).
\item ΔCPM S4 − S3 per condition.
\item The above by native kernel class.
\item Candidates per scan.
\item Measured TCIA noise correlation between adjacent slices (descriptive; checks the z-independence assumption).
\end{itemize}

\subsubsection*{Interpretation rule (fixed now; the paper's wording follows it)}
\begin{itemize}
\item \textbf{Supported}: F4 significant for both models AND the equivalence check holds for both. Texture at equal level explains most of the difference, and stationary NPS-matched noise suffices to simulate dose.
\item \textbf{Partly supported}: F4 significant for at least one model, but the equivalence check fails. Texture matters, but so does the non-stationary or streak structure of projection noise; NPS-matched stationary noise is not a full substitute.
\item \textbf{Not supported}: F4 significant for neither model. At equal level and path, texture does not explain the difference; the draft's explanation must be withdrawn.
\end{itemize}

The result is reported whichever way it goes.

\subsubsection*{Changes after this point}
(none yet)

